\clearpage
\section{세 고전 불가능 문제}
\label{sec:construction-classical-impossibilities}

\begin{learninggoal}
정육면체 배적, 각의 삼등분과 원적문제를 체확장 차수와 초월성으로 판정한다. ``특정한 경우의 가능성''과 ``임의의 입력에 대한 보편 작도''를 구분한다.
\end{learninggoal}

고대 그리스의 세 유명한 문제는 모두 어떤 길이나 각을 자와 컴퍼스로 만들 수 있는지 묻는다. 갈루아 이론의 언어로 바꾸면, 목표 좌표가 $\mathcal C_0$에 속하는지를 묻는 문제가 된다.

\subsection{정육면체의 배적}

모서리 길이가 $1$인 정육면체의 부피는 $1$이다. 부피가 두 배인 정육면체의 모서리 길이를 $x$라 하면
\[
  x^3=2,
  \qquad
  x=\sqrt[3]{2}
\]
이다.

\begin{theorem}[정육면체 배적의 불가능성]
자와 컴퍼스만으로 단위 정육면체의 부피를 두 배로 만드는 정육면체를 작도할 수 없다.
\end{theorem}

\begin{proof}
다항식 $X^3-2$는 소수 $2$에 대한 아이젠슈타인 판정법으로 $\Q$ 위에서 기약이다. 따라서
\[
  [\Q(\sqrt[3]2):\Q]=3.
\]
작도가능한 대수적 수의 차수는 $2$의 거듭제곱이어야 하는데 $3$은 그렇지 않다. 그러므로 $\sqrt[3]2$는 작도가능하지 않다.
\end{proof}

이 증명은 단순히 ``세제곱근은 컴퍼스로 못 그린다''고 말하지 않는다. 예를 들어 $\sqrt[3]8=2$는 당연히 작도가능하다. 핵심은 특정한 목표수 $\sqrt[3]2$의 최소다항식 차수가 $3$이라는 점이다.

\subsection{임의의 각의 삼등분}

어떤 각들은 쉽게 삼등분된다. 예를 들어 직각 $90^\circ$는 $30^\circ$씩 나눌 수 있다. 고전 문제는 \emph{주어진 임의의 각을 항상} 자와 컴퍼스로 삼등분하는 절차가 존재하는지를 묻는다. 이를 부정하려면 한 각만 반례로 제시하면 충분하다.

\begin{theorem}[$60^\circ$의 삼등분 불가능성]
자와 컴퍼스로 $60^\circ$를 $20^\circ$씩 삼등분할 수 없다. 따라서 임의의 각을 삼등분하는 보편적인 자와 컴퍼스 작도는 존재하지 않는다.
\end{theorem}

\begin{proof}
$20^\circ$가 작도가능하다고 가정하자. 그러면 단위원 위의 점
\[
  e^{i\pi/9}
\]
과 그 실수부 $\cos(\pi/9)$도 작도가능하다. 편의를 위해
\[
  y=2\cos(\pi/9)
\]
라 두자. 삼중각 공식
\[
  2\cos(3\theta)=(2\cos\theta)^3-3(2\cos\theta)
\]
에 $\theta=\pi/9$를 대입하면
\[
  1=y^3-3y
\]
이므로
\[
  y^3-3y-1=0.
\]
유리근 정리에 따른 후보 $\pm1$은 근이 아니므로 $X^3-3X-1$은 $\Q$ 위에서 기약이다. 따라서
\[
  [\Q(y):\Q]=3,
\]
이는 작도가능성의 차수 조건과 모순이다.
\end{proof}

\begin{remark}
같은 계산은 정구각형이 작도 불가능하다는 사실과 연결된다. 정구각형의 중심각은 $40^\circ$이고, 이를 이등분하면 $20^\circ$가 된다. 따라서 정구각형을 작도할 수 있었다면 $60^\circ$의 삼등분도 가능했을 것이다.
\end{remark}

\subsection{원의 정사각화}

반지름 $1$인 원의 넓이는 $\pi$이다. 같은 넓이를 갖는 정사각형의 한 변의 길이는
\[
  \sqrt\pi
\]
이다.

\begin{theorem}[원적문제의 불가능성]
자와 컴퍼스만으로 주어진 원과 같은 넓이의 정사각형을 작도할 수 없다.
\end{theorem}

\begin{proof}
표준 작도가능체 $\mathcal C_0$의 모든 원소는 $\Q$ 위에서 대수적이다. 그러나 린데만의 정리에 의해 $\pi$는 $\Q$ 위에서 초월적이다. 만약 $\sqrt\pi$가 대수적이라면 그 제곱 $\pi$도 대수적이어야 하므로 $\sqrt\pi$ 역시 초월적이다. 따라서 $\sqrt\pi\notin\mathcal C_0$이다.
\end{proof}

린데만의 초월성 정리는 이 교재의 범위를 넘어서는 깊은 결과이다. 여기서 중요한 것은 작도이론이 초월성 정리와 결합될 때 원적문제를 즉시 해결한다는 구조이다.

\begin{center}
\renewcommand{\arraystretch}{1.25}
\begin{tabular}{llll}
\toprule
문제 & 목표량 & 장애 & 결론 \\
\midrule
정육면체 배적 & $\sqrt[3]2$ & 최소다항식 차수 $3$ & 불가능 \\
$60^\circ$ 삼등분 & $2\cos20^\circ$ & 최소다항식 차수 $3$ & 불가능 \\
원의 정사각화 & $\sqrt\pi$ & 초월수 & 불가능 \\
\bottomrule
\end{tabular}
\end{center}

\begin{pitfall}
\begin{itemize}
  \item ``각의 삼등분이 불가능하다''는 말은 모든 특정 각이 불가능하다는 뜻이 아니다. 임의의 입력각을 처리하는 보편 작도가 없다는 뜻이다.
  \item 매우 정확한 근사작도와 정확한 유한 작도는 다르다. 이 장은 정확한 작도만을 다룬다.
  \item 작도 불가능성을 보일 때 차수가 $2$의 거듭제곱이 아님을 확인하는 방법은 충분하지만, 차수가 $2$의 거듭제곱이라고 해서 자동으로 작도가능한 것은 아니다.
\end{itemize}
\end{pitfall}

\begin{checkpoint}
\begin{enumerate}[label=(\alph*)]
  \item $90^\circ$가 삼등분 가능하다는 사실은 위 정리와 왜 모순되지 않는가?
  \item $\sqrt[6]2$의 차수를 구하고 작도가능성을 판정하라.
  \item 원적문제의 증명에서 체확장론 외에 필요한 외부 정리는 무엇인가?
\end{enumerate}
\end{checkpoint}
